\documentclass[letterpaper]{article}
\usepackage[submission]{aaai2027}
\usepackage[hyphens]{url}
\usepackage{graphicx}
\urlstyle{rm}
\def\UrlFont{\rm}
\usepackage{natbib}
\usepackage{caption}
\frenchspacing
\usepackage{algorithm}
\usepackage{algorithmic}
\usepackage{booktabs}
\usepackage{amsmath}
\usepackage{amssymb}
\usepackage{amsfonts}
\usepackage{microtype}
\usepackage{amsthm}
\newtheorem{theorem}{Theorem}
\newtheorem{proposition}{Proposition}
\pdfinfo{
/TemplateVersion (2027.1)
}

\setcounter{secnumdepth}{0}

\title{Supplementary Material:\\
TEEG-INCR: Millisecond-Scale Incremental Updates for\\
Public-Transit Earliest-Arrival Routing at City Scale}

\author{Anonymous Authors}
\affiliations{Anonymous Affiliations}

\begin{document}
\maketitle

\section*{Technical Appendix}
\setcounter{section}{0}
\setcounter{theorem}{0}
\setcounter{table}{0}
\renewcommand{\thesection}{\Alph{section}}
\renewcommand{\thetable}{S\arabic{table}}

This supplementary material accompanies the main paper. References to
``Proposition~1'', ``Section~6.4'', ``Table~1'', etc. point to the main
paper unless prefixed with ``S'' (e.g.\ Table~S3).

\section{Proof of Proposition 1}
\label{app:proof}

We restate the proposition for clarity, then prove it.

\begin{proposition}[Empirical EA-label domination, restated]
Fix a TEEG instance built from a public-transit network with stop set
$S$, connection set $C$, footpath weights $W$, and a query
$(s_\text{src}, s_\text{tgt}, t_\text{start})$. Let $\tau_\textsc{csa}[s]$
be the EA label produced by single-pass CSA (Dibbelt et al.\ 2018,
Algorithm~1) and let $\tau^S_\textsc{teeg}[s] := \min_{v: \text{stop}(v) = s}
\tau^V_\textsc{teeg}[v]$ be the EA label produced by TEEG-ALT. Then for
every $s \in S$, $\tau^S_\textsc{teeg}[s] \leq \tau_\textsc{csa}[s]$.
The inequality is \emph{strict} ($\tau^S_\textsc{teeg}[s] <
\tau_\textsc{csa}[s]$) for some $s \in S$ if any of:
\begin{enumerate}
\item[(i)] $\text{stop}^\text{dep}_c \neq s_\text{src}$ and
$W[s_\text{src}, \text{stop}^\text{dep}_c] / v_\text{walk} \leq
t^\text{dep}_c - t_\text{start} \leq W_\text{max}/v_\text{walk}$;
\item[(ii)] $\text{stop}^\text{dep}_{c'} = s_\text{src}$ and
$t^\text{dep}_{c'} \leq t^\text{dep}_c$, with CSA scanning
connections in non-decreasing $t^\text{dep}$ order single-pass; and
\item[(iii)] there exists a target $s^* \in S$ such that the EA at
$s^*$ via $c$ is strictly earlier than the EA at $s^*$ via $c'$;
\end{enumerate}
\emph{or} (iv) the two-hop walk-chaining configuration:
there exist stops $a, b \in S$ and a connection $c''$ with
$\text{stop}^\text{dep}_{c''} = b$ such that
$(W[s_\text{src}, a] + W[a, b]) / v_\text{walk} \leq
t^\text{dep}_{c''} - t_\text{start}$
yet $W[s_\text{src}, b] / v_\text{walk} > t^\text{dep}_{c''} -
t_\text{start}$ (i.e., a two-hop walk reaches $b$ in time but no
single walk does), and the EA at some $s^* \in S$ via $c''$ is
strictly earlier than via any connection from a one-hop walk-reachable
stop.
\end{proposition}

\paragraph{Inclusion ($\tau^S_\textsc{teeg} \leq \tau_\textsc{csa}$ pointwise).}
We prove this by simulating CSA's single-pass connection scan with
TEEG-ALT's event relaxation pass. Let $R_k$ denote the $k$-th
connection that CSA relaxes (in non-decreasing $t^\text{dep}$ order).
We prove by induction on $k$ that after TEEG-ALT processes all heap
entries with key $\leq t^\text{dep}_{R_k}$, we have
$\tau^V_\textsc{teeg}[v^\text{arr}_{R_k}] \leq \tau_\textsc{csa}[\text{stop}^\text{arr}_{R_k}]$.
\emph{Base case ($k=0$):} CSA's source-walk initialisation
(Alg.~1, lines 3--8) writes $\tau[\text{stop}_x] \leftarrow t_\text{start}
+ W[s_\text{src}, \text{stop}_x]/v_\text{walk}$ for every walk-reachable
$\text{stop}_x$. TEEG-ALT's source-walk loop (Algorithm~\ref{alg:teeg-alt},
lines 5--10) writes $\tau^V[v^\text{dep}_c] \leftarrow \min(\tau^V[v^\text{dep}_c], t^\text{dep}_c)$
for every walk-reachable departure event $c$ with
$t^\text{dep}_c \geq \tau_\text{walk}$, simulating each CSA initialisation
write through the dwell-edge $\text{stop}_x \rightsquigarrow v^\text{arr}_c$.
\emph{Inductive step:} the $k$-th CSA relaxation triggers one of:
(a) ride: $\tau[\text{stop}^\text{arr}_{R_k}] \leftarrow t^\text{arr}_{R_k}$
if $\tau[\text{stop}^\text{dep}_{R_k}] \leq t^\text{dep}_{R_k}$;
(b) trip-continue (CSA's implicit ``stay on vehicle'' across consecutive
connections of the same trip);
(c) walk closure from $\text{stop}^\text{arr}_{R_k}$.
All three have direct TEEG analogues with weight-preserving edges:
(a) the RIDE edge $v^\text{dep}_{R_k} \to v^\text{arr}_{R_k}$ with
weight $t^\text{arr}_{R_k} - t^\text{dep}_{R_k}$;
(b) the DWELL edge $v^\text{arr}_{R_k} \to v^\text{dep}_{R_{k+1}}$ for
$R_k, R_{k+1}$ consecutive on the same trip;
(c) WALK edges from $v^\text{arr}_{R_k}$ to each $v^\text{dep}_{c'}$
with $\text{stop}^\text{dep}_{c'}$ walk-reachable from
$\text{stop}^\text{arr}_{R_k}$ \emph{and} TRANSFER edges for same-stop
boarding with $t^\text{dep}_{c'} - t^\text{arr}_{R_k} \geq
\tau_\text{transfer}$. By the induction hypothesis applied to the source
event $v^\text{dep}_{R_k}$, the LHS of each TEEG relaxation is $\leq$
the LHS of the corresponding CSA relaxation. TEEG-ALT is a
Dijkstra-style label-correcting algorithm with a consistent monotone
heuristic (Lemma~A.0) on a graph with non-negative edge weights;
standard A$^*$/Dijkstra optimality then guarantees
$\tau^V[v]$ converges to the true shortest-path EA at $v$ before $v$
is popped from the heap. The pointwise inequality follows by taking
the minimum over all events at each stop. For stops $s$ with no
outgoing connection in $[t_\text{start}, t_\text{end}]$, we extend
$\tau^S_\textsc{teeg}[s] := t_\text{start} + W[s_\text{src},
s]/v_\text{walk}$ if $s$ is walk-reachable from $s_\text{src}$ within
$W_\text{max}$, else $\infty$; this matches CSA's source-walk init
and preserves the inequality. \qed

\paragraph{Strict inequality under (i)--(iv).} Construct the witness.
Under (i)--(iii), CSA's single-pass discipline (Alg.~1 processes each
connection exactly once in $t^\text{dep}$ order) commits to
$c'$ first because $t^\text{dep}_{c'} \leq t^\text{dep}_c$. Downstream
relaxations from $c'$ raise $\tau$ at intermediate stops, suppressing
the walk-relaxation that would propagate from $c$. TEEG-ALT, by
contrast, materialises $c$ and $c'$ as two independent event vertices
with two independent source-edges, both initialised in lines 5--10.
The heap-driven relaxation pass then propagates the strictly-better
arrival at $s^*$ from $c$, yielding $\tau^S_\textsc{teeg}[s^*] <
\tau_\textsc{csa}[s^*]$.
Under (iv), the second-hop walk is never relaxed by CSA at all
(closure cap), so $\tau_\textsc{csa}[b] = \infty$ while TEEG materialises
the two-hop walk via its WALK-of-WALK edges. \qed

\paragraph{ALT heuristic admissibility (Lemma A.0).}
The landmark distances $d_l(v)$ are computed as event-graph
shortest-path distances to landmark $l$. Define
$h[v] := \max_l \max(0, d_l(v) - \max_{u : \text{stop}(u) = s_\text{tgt}} d_l(u))$,
matching Algorithm~\ref{alg:teeg-alt} line~4. (\emph{Admissibility.})
Let $u^* := \arg\min_{u : \text{stop}(u) = s_\text{tgt}} \text{dist}(v, u)$
be the closest target-stop event to $v$ in the event graph. For each
landmark~$l$, the reverse triangle inequality gives $d_l(v) - d_l(u^*)
\leq \text{dist}(v, u^*)$. Taking $\max_u d_l(u)$ over target-stop
events $u$ can only enlarge $d_l(u^*)$ (a single instance of $u$),
hence $d_l(v) - \max_u d_l(u) \leq d_l(v) - d_l(u^*) \leq \text{dist}(v, u^*)$.
Taking the outer $\max_l \max(0, \cdot)$ preserves this lower bound:
$h[v] \leq \text{dist}(v, u^*) = \min_u \text{dist}(v, u)$, which is
the optimal cost to any target-stop event from $v$. Hence $h$ is
admissible. (\emph{Consistency.}) For every edge $(v, u) \in E$,
$d_l$ is itself a shortest-path distance so $d_l(v) - d_l(u) \leq w(v, u)$;
this implies $h(v) - h(u) \leq w(v, u)$ since the outer
$\max$/$\max(0, \cdot)$ are 1-Lipschitz. Hence $h$ is consistent, and
Algorithm~\ref{alg:teeg-alt}'s deferred-termination / goal-pruning is
sound. \qed

\paragraph{Lemma A.0' --- ALT admissibility survives TEEG-INCR updates.}
TEEG-INCR's query path (Algorithm~\ref{alg:teeg-incr-query}) reuses
the precomputed landmark distance vector $\{d_l\}$ from the static
substrate without recomputation per edit. We claim pre-update
$\{d_l\}$ remains an admissible (and consistent) lower bound on
post-update EA distance.

\emph{Proof.} Let $G$ be the pre-update event graph and $G'$ be
the post-update graph after a sequence of \textsc{cancel} /
\textsc{delay} / \textsc{add} primitives. We show
$\text{dist}_{G'}(v, u^*) \geq \text{dist}_G(v, u^*) - c$ for some
non-negative constant $c$ derivable from the primitives:

\begin{itemize}
\item \textsc{cancel}$(e)$ removes connection event vertices and
their incident DWELL / TRANSFER / WALK\_AND\_BOARD edges (or
tombstones them in the overlay). Removing edges from a graph
\emph{cannot decrease} shortest-path distance:
$\text{dist}_{G \setminus E_e}(v, u^*) \geq \text{dist}_G(v, u^*)$.
Therefore $d_l^G(v) \leq d_l^{G'}(v) \leq \text{dist}_{G'}(v, u^*)$
(taking $u^* = $ landmark $l$).
\item \textsc{delay}$(e, \Delta)$ tombstones the original arrival
event and inserts a delayed-arrival event $\Delta$ later. The
delayed event-vertex has all outgoing DWELL / TRANSFER edges of
the original but starts $\Delta$ later, so any path that uses
$e$'s arrival as an intermediate vertex is $\geq \Delta$ longer.
Other paths are unaffected. Therefore
$\text{dist}_{G'}(v, u^*) \geq \text{dist}_G(v, u^*)$ (paths
through $e$) or $= \text{dist}_G(v, u^*)$ (paths not through $e$).
\item \textsc{add}$(e)$ inserts new event vertices not in $\{d_l\}$.
Querying $d_l(v)$ for $v$ a newly-added event falls back to the
conservative stop-level bound
$d_l(v) := \min_{u' : \text{stop}(u') = \text{stop}(v),~u' \in V_G} d_l(u')$,
which is a lower bound on the new vertex's true distance to $l$
because adding edges to $G$ does not invalidate $\{d_l\}$ as a
lower bound for existing vertices, and the stop-level minimum is
admissible by the same triangle-inequality argument as in
Lemma A.0.
\end{itemize}

Combining: $\{d_l^G\}$ is an admissible (and consistent, since
the same 1-Lipschitz argument applies) lower bound on post-update
EA distance. ALT goal-pruning remains correct; pruning may become
slightly looser (more nodes relaxed than strictly necessary) but
not incorrect. \qed

The 10{,}800-comparison oracle (Appendix~D) confirms zero
correctness regressions attributable to stale landmarks across all
36 Cycle~10 cells. A periodic landmark refresh (cost:
$8\,\mathrm{s}$ for London per Table~4) restores tight pruning if
drift accumulates over very long update sessions.

\paragraph{Empirical instantiation.} We verified all 17 (of 1{,}500)
Full-London queries where TEEG-ALT and single-pass CSA disagree by
edge-by-edge timetable feasibility (every RIDE edge against the GTFS
connection set, every WALK edge against $W/v_\text{walk}$, every
TRANSFER edge against the stop's minimum-transfer-time). In 4/17,
conditions (i)--(iii) (single-hop transit witness) apply with
$|\Delta|=4\text{--}21$\,min and \texttt{n\_csa\_missed\_on\_teeg\_path}\,$>0$,
with $c$ a Tube connection a 4--9 minute walk from $s_\text{src}$ and
$c'$ a bus departing from $s_\text{src}$. The remaining 13/17 fall under
condition (iv) (two-hop walk-chain), $|\Delta|=1\text{--}6$\,min,
\texttt{n\_csa\_missed\_on\_teeg\_path}\,$=0$: $s_\text{src} \to a \to b$
opens an earlier Tube connection at $b$ that the single-hop
walk-closure from $s_\text{src}$ does not reach.

\paragraph{Verification artefacts.} The 17 per-query traces (with
edge-by-edge feasibility audit, CSA $\tau$ trace, and TEEG event-path
reconstruction) are at
\texttt{Experiments/results/cycle\_9\_paper/RESIDUAL\_17\_DEEP\_DIVE/\\q\{102,430,431,634,827,930,1384,1432,1497,\\1640,1800,1811,1858,1893,2038,2415,2759\}.md}
in the supplementary code release. The feasibility script
\texttt{Experiments/cycle\_runner/deep\_dive\_17\_residuals.py} produces
this audit deterministically (seed \texttt{numpy.random.seed(20260101)},
TEEG build $\sim 12$\,min wall, per-query analysis $\sim 5$\,s).

All 17 cases were independently hand-audited: 5 (q102, q1640, q431,
q1384, q1893) by direct trace-from-GTFS, and the remaining 12
(q430, q634, q827, q930, q1432, q1497, q1800, q1811, q1858, q2038,
q2415, q2759) by human cross-check of the deterministic script's
emitted per-edge \texttt{OK} annotations against \texttt{b.connections}
+ \texttt{transfers\_idx} and per-stop $\tau$ comparison. \textbf{All
17 audits returned TEEG\_CORRECT with high confidence; zero
TEEG\_BUG, zero inconclusive.} The 17 residuals split into two
known one-pass-CSA failure-mode subtypes (Dibbelt 2018, Sec.~3.4):
\emph{src-walk-suppresses} (4 cases, $\Delta\!=\!4\text{--}9$\,min)
and \emph{walk-from-walk} (13 cases, $\Delta\!=\!1\text{--}6$\,min,
$n_\text{csa\_missed\_on\_teeg\_path}\!=\!0$). Per-query verdicts at
\texttt{handtrace\_12\_residuals/q*\_handtrace.md}; summary at
\texttt{handtrace\_12\_residuals/summary.md}.

\paragraph{Cross-city generalisation.} The structural argument
generalises to Berlin and NYC by the same construction (conditions
(i)--(iv) depend only on TEEG/CSA semantics, not on London-specific
data). The 17 empirical witnesses are London-only; on Berlin and NYC
we observe similar TEEG vs single-pass-CSA disagreement rates
(\textasciitilde 0.8\% Berlin, \textasciitilde 0.4\% NYC, both below
London's 1.1\%), consistent with the same failure modes. Witness-level
audit on those cities is a Phase II item.

\paragraph{Independent oracle: multi-pass CSA, full 3{,}000-query sweep.}
To independently validate that single-pass CSA has a real algorithmic
blind spot at Full London scale (and not just on a 17-query subset
TEEG-ALT happened to disagree on), we implemented Dibbelt et al.\
\cite{dibbelt2018csa} Algorithm~2 (multi-pass CSA): walk-relax
unconditionally from every usable connection's arrival stop and
iterate the connection scan to fixpoint, resetting the per-trip
boarding flag each pass. Walks are $F_{15}$-closed (transitive
closure restricted to per-leg paths of duration $\leq 15$\,min).
We ran this on \emph{all} 3{,}000 Full London queries (not just the
17 TEEG-ALT residuals).

\textbf{Results on the 3{,}000-query sweep.} With
$\Delta \!:=\! \text{EA}_\text{multipass}\!-\!\text{EA}_\text{singlepass}$
(so $\Delta\!<\!0$ means multipass returns an earlier arrival),
multi-pass CSA Algorithm~2 returns a strictly earlier EA than
single-pass Algorithm~1 on \textbf{185 queries (6.17\%)}, with
median $\Delta\!=\!-8$\,min, range $-24$ to $-1$\,min, and p95 of
the multi-pass per-query wall-time of $0.19$\,s. The 17 prior
TEEG-ALT-vs-singlepass residuals are a verified subset of these
185 (consistency check passed; q\_id=1497 has $\Delta\!=\!-2$\,min
in this sweep, confirming it belongs in the 185).
Mode breakdown of the 185 (percentages rounded; may not sum to
exactly 100\%):
\begin{itemize}
\item large source-walk-suppresses ($|\Delta|\!=\!6\text{--}15$\,min): 110 (59.5\%)
\item medium transfer-chain ($|\Delta|\!=\!3\text{--}5$\,min): 36 (19.5\%)
\item small walk-from-walk ($|\Delta|\!\leq\!2$\,min): 30 (16.2\%)
\item huge alternative-route ($|\Delta|\!>\!15$\,min): 9 (4.9\%)
\end{itemize}

\textbf{Results on the 17 TEEG-ALT residuals specifically.} Within
the 17 queries where TEEG-ALT and single-pass CSA disagree,
multi-pass CSA agrees with TEEG-ALT to the minute on \textbf{16/17}.
On the remaining query (\texttt{q\_id\!=\!1497}), multi-pass CSA
finds an arrival 1\,minute earlier than TEEG-ALT (642 vs 643\,min)
via a separate Ivydale-Road walk-and-board chain that TEEG-ALT's
goal-pruning misses. This confirms the gap diagnosed in
\S\ref{app:proof} is the textbook
``src-walk-suppresses-post-arrival-walk-relax'' blind spot
(Dibbelt et al.\ \cite{dibbelt2018csa}), closed by the multi-pass
fixpoint variant. Per-query results for the 17-residual
investigation in
\texttt{multipass\_csa\_oracle/results\_17.csv}; the full
3{,}000-query sweep in
\texttt{multipass\_3000\_full/results\_3000.csv} +
\texttt{multipass\_3000\_full/residual\_classification.csv} +
\texttt{multipass\_3000\_full/summary.json}; implementation in
\texttt{Route\_Planner/reference\_oracle/csa\_multipass.py}.

\paragraph{Note on the original Reference-CSA${}^*$ oracle.}
Our cross-impl single-pass Reference-CSA${}^*$ (Algorithm~1)
returns the \emph{same} EA values as production CSA on all 17
queries. This is expected: both implement the same single-pass
algorithm and inherit the same blind spot. Reference-CSA${}^*$'s
role is detecting \emph{implementation} bugs in our production CSA
(it caught a dwell-row-splitting bug during Cycle~9), not
certifying single-pass CSA as Pareto optimal. The multi-pass
oracle above is the structurally-independent algorithmic check.

\section{Algorithm Pseudocode}
\label{app:alg}

\begin{algorithm}[H]
\caption{TEEG-ALT static EA query}
\label{alg:teeg-alt}
\begin{algorithmic}[1]
\REQUIRE TEEG $G=(V,E)$, source stop $s_\text{src}$, target stop $s_\text{tgt}$, start time $t_\text{start}$, event-level landmark distances $\{d_l(\cdot)\}$
\STATE \textbf{if} $s_\text{src} = s_\text{tgt}$ \textbf{then return} $t_\text{start}$
\STATE $\text{best} \gets t_\text{start} + W[s_\text{src}, s_\text{tgt}]/v_\text{walk}$ if walk-reachable else $\infty$
\STATE $\tau^V[v] \gets \infty$ for all $v \in V$
\STATE $h[v] \gets \max_l \max(0, d_l(v) - \max_{u: \text{stop}(u)=s_\text{tgt}} d_l(u))$\quad{\small // admissible ALT lower bound}
\FORALL{$v^\text{dep}_c$ with $\text{stop}^\text{dep}_c$ walk-reachable from $s_\text{src}$}
  \STATE $\tau_\text{walk} \gets t_\text{start} + W[s_\text{src}, \text{stop}^\text{dep}_c]/v_\text{walk}$
  \IF{$\tau_\text{walk} \leq t^\text{dep}_c$}
    \STATE $\tau^V[v^\text{dep}_c] \gets \min(\tau^V[v^\text{dep}_c], t^\text{dep}_c)$
  \ENDIF
\ENDFOR
\STATE \texttt{heap} $\gets$ min-heap initialised with $\{(\tau^V[v] + h[v], v) : \tau^V[v] < \infty\}$
\WHILE{\texttt{heap} not empty}
  \STATE $(k, v) \gets$ \texttt{heap.pop()}
  \IF{$k > \tau^V[v] + h[v]$} \STATE \textbf{continue}\quad{\small // stale entry guard} \ENDIF
  \IF{$\text{stop}(v) = s_\text{tgt}$}
    \STATE $\text{best} \gets \min(\text{best}, \tau^V[v])$; \textbf{continue}\quad{\small // deferred termination}
  \ENDIF
  \IF{$k > \text{best}$}
    \RETURN $\text{best}$\quad{\small // ALT goal pruning vs current best}
  \ENDIF
  \FORALL{$(v, u) \in E$}
    \STATE $\tau' \gets \tau^V[v] + w(v, u)$
    \IF{$\tau' < \tau^V[u]$}
      \STATE $\tau^V[u] \gets \tau'$
      \STATE \texttt{heap.push}$((\tau' + h[u], u))$
    \ENDIF
  \ENDFOR
\ENDWHILE
\RETURN $\text{best}$
\end{algorithmic}
\end{algorithm}

\begin{algorithm}[H]
\caption{TEEG-INCR \textsc{cancel}/\textsc{delay}/\textsc{add}/\textsc{query}/\textsc{compact}}
\label{alg:teeg-incr}
\begin{algorithmic}[1]
\REQUIRE TEEG state $\Sigma$ (CSR $G$, tombstone bitvec $T$, overlay $\Omega$, dirty list $D$)
\STATE \textbf{procedure} \textsc{cancel}($c$):
\STATE \quad $T[v^\text{dep}_c] \gets 1$; $T[v^\text{arr}_c] \gets 1$
\STATE \quad \textbf{for} $u \in \text{out-nbrs}(v^\text{dep}_c) \cup \text{out-nbrs}(v^\text{arr}_c)$ \textbf{do} $D.\textsc{push}(u)$
\STATE \textbf{procedure} \textsc{delay}($c$, $\Delta$):
\STATE \quad $T[v^\text{dep}_c] \gets 1$; $T[v^\text{arr}_c] \gets 1$
\STATE \quad allocate ${v'}^\text{dep}_c, {v'}^\text{arr}_c$ in $\Omega$ with times shifted by $\Delta$
\STATE \quad insert outgoing dwell/transfer/walk targets for ${v'}^\text{arr}_c$ into $\Omega$
\STATE \quad \textsc{build-in-incidence}$({v'}^\text{dep}_c)$: \quad{\small // scan in-neighbours of $v^\text{dep}_c$ in $(G\!\setminus\!T)$ and redirect each $(u, v^\text{dep}_c)$ to $(u, {v'}^\text{dep}_c)$ in $\Omega$ when the time constraint $t(u) + w(u, {v'}^\text{dep}_c) \leq t^\text{dep}_c + \Delta$ still holds}
\STATE \quad \textbf{for} $u \in \text{out-nbrs}({v'}^\text{arr}_c) \cup \text{in-nbrs}({v'}^\text{dep}_c)$ \textbf{do} $D.\textsc{push}(u)$
\STATE \textbf{procedure} \textsc{add}($c$):
\STATE \quad allocate $v^\text{dep}_c, v^\text{arr}_c$ in $\Omega$ and insert outgoing incidences via build-time helper
\STATE \quad \textsc{build-in-incidence}$(v^\text{dep}_c)$
\STATE \quad \textbf{for} $u \in \text{out-nbrs}(v^\text{arr}_c) \cup \text{in-nbrs}(v^\text{dep}_c)$ \textbf{do} $D.\textsc{push}(u)$
\STATE \textbf{procedure} \textsc{query}($s_\text{src}, s_\text{tgt}, t_\text{start}$):
\STATE \quad \textbf{while} $D$ not empty: $v \gets D.\textsc{pop}()$; recompute $\tau^V[v]$ over in-edges of $v$ in $(G \setminus T) \cup \Omega$
\STATE \quad \textbf{return} \textsc{TEEG-ALT}($(G \setminus T) \cup \Omega, s_\text{src}, s_\text{tgt}, t_\text{start}, \{d_l\}$)
\STATE \textbf{procedure} \textsc{compact}():
\STATE \quad \textbf{if} $\textsc{popcount}(T)/|V| > 0.1$ \textbf{then} rebuild CSR $G$ skipping tombstones, fold $\Omega$ in, clear $T, \Omega, D$
\end{algorithmic}
\end{algorithm}

\begin{algorithm}[H]
\caption{\textsc{build-in-incidence} helper used by Algorithm~2}
\label{alg:build-in-incidence}
\begin{algorithmic}[1]
\REQUIRE TEEG state $\Sigma$, target event $v$ to install in-edges for
\STATE $S \gets \text{stop}(v)$, $t \gets \text{time}(v)$
\FORALL{event $u$ with $\text{stop}(u)$ walk-/transfer-reachable to $S$}
  \STATE $w \gets$ minimum legal traversal time from $u$ to $v$ (transfer / walk / dwell, per edge type)
  \IF{$u \notin T$ \textbf{and} $\text{time}(u) + w \leq t$}
    \STATE insert $(u, v, w)$ into $\Omega$\quad{\small // overlay carries new in-edge}
  \ENDIF
\ENDFOR
\end{algorithmic}
\end{algorithm}

\section{Cross-implementation deep-dive}
\label{app:kit}

The main paper's Table~3 (KIT comparison) reports the agreement of our
implementations against the publicly available KIT ULTRA
implementation~\cite{ulra-kit} on 3{,}000 Full London queries (47
dropped due to KIT's maximum-connected-component reduction, leaving
2{,}953 mutually-runnable ODs). The two pipelines differ in two
independent dimensions: \emph{scan / data-pipeline} (CSA scan
implementation + service-day filter + connection encoding + stop set)
and \emph{walking-transfer graph} (network-aware corrector vs
haversine $\leq\!900$\,s).

\paragraph{4-way controlled ablation.} To attribute the 55.8\,pp gap
(100\% $-$ 44.2\%) cleanly, we ran a $2\times 2$ ablation on a 500-OD
subset, crossing \{KIT scan, our scan\} with \{KIT walks, our walks\}.
Cell~2 (KIT C++ scan + our walking graph) was produced by patching
KIT's source with a new \texttt{loadTransfersFromCsv} runnable
(73 LoC; \texttt{cell2\_patch.diff} in the supplementary code release)
that overwrites \texttt{Intermediate::TransferGraph} with an external
CSV before \texttt{intermediateToCSA}.

\begin{table}[H]
\centering
\caption{Cross-pipeline pairwise agreement on 500 Full-London OD pairs
($\pm 1$ min tolerance). Cell~1: KIT scan + KIT walks (vanilla).
Cell~2: KIT scan + our walks. Cell~3: our scan + KIT walks. Cell~4:
our scan + our walks (vanilla).}
\label{tab:kit-4way}
\small
\begin{tabular}{lrrrr}
\toprule
              & cell~1 & cell~2 & cell~3 & cell~4 \\
\midrule
cell~1 & 100.0\% & 79.4\% & 42.8\% & 44.2\% \\
cell~2 &  79.4\% & 100.0\% & 39.4\% & 40.0\% \\
cell~3 &  42.8\% &  39.4\% & 100.0\% & 88.2\% \\
cell~4 &  44.2\% &  40.0\% &  88.2\% & 100.0\% \\
\bottomrule
\end{tabular}
\end{table}

\paragraph{Decomposition (corroborated from both diagonals).}
\begin{itemize}
\item \textbf{Walking-graph contribution: $12\text{--}21$\,pp.}
Cell~3 vs Cell~4 (our scan, walks differ) gives $11.8$\,pp;
Cell~1 vs Cell~2 (KIT scan, walks differ) gives $20.6$\,pp; average
$\sim\!16$\,pp.
\item \textbf{Scan / data-pipeline contribution: $57\text{--}60$\,pp.}
Cell~1 vs Cell~3 (KIT walks, scan differs) gives $57.2$\,pp;
Cell~2 vs Cell~4 (our walks, scan differs) gives $60.0$\,pp; average
$\sim\!58$\,pp.
\end{itemize}
Both diagonals corroborate the same decomposition.

\paragraph{Direction.} The signed median $\Delta$ (our scan minus
KIT) is $+10.0$\,min with KIT walks and $+12.5$\,min with our walks,
stable across both walking-graph configurations. Our oracle is
uniformly \emph{earlier}, not randomly different --- an
information-set inequality, not a scoring error.

\paragraph{Per-suspect attribution of the 55.8\,pp scan-side gap.}
Re-running KIT under three patched configurations (S1 service-day
filter, S2 int16 deci-seconds time quantization, S3 MCC stop-set)
would each require intrusive C++ patches to KIT's build tree. As a
faster alternative, we \emph{partition} the existing 3{,}000-OD
KIT-vs-oracle disagreement data by the empirical signature each
suspect leaves. Using
$\Delta \!:=\! \text{EA}_\text{oracle}\!-\!\text{EA}_\text{KIT}$
(so $\Delta\!<\!0$ means KIT thinks the trip is slower than the
oracle), we classify every disagreement into one of three named
suspects (S1, S3) plus two residual buckets (S5a, S5b); S2 is
\emph{not} attributable from external data alone because int16-deci-second
quantization at KIT's encoding stage leaves a sub-second signature
that gets folded into the bulk residual S5. We use priority
$S3>S1>(S5\text{ by sign of }\Delta)$ to ensure each disagreement
is counted once:
\begin{itemize}
\item \textbf{S3 MCC stop-set} (\texttt{kit\_dropped} $=$ True \textbf{or}
\texttt{agree\_csa} $=$ \texttt{ORACLE\_FINDS\_KIT\_BLIND}): 125 ODs,
$\mathbf{4.2\,\text{pp}}$ of the $55.8$\,pp.
\item \textbf{S1 service-day filter} ($|\Delta|\!>\!60$\,min cliff,
diagnostic of trips that span midnight): 179 ODs, $\mathbf{6.0\,\text{pp}}$.
\item \textbf{S5a KIT-slower-than-oracle} (signed $\Delta\!<\!0$,
remaining after S1+S3): 1{,}331 ODs, $\mathbf{44.6\,\text{pp}}$.
This is the dominant residual; consistent with a window-cutoff,
restrictive route-filter, or stricter MTT inside KIT's
\texttt{intermediate.binary} that does not trigger the 60-min cliff.
\item \textbf{S5b KIT-faster-than-oracle} (signed $\Delta\!>\!0$,
remaining): 31 ODs, $\mathbf{1.0\,\text{pp}}$. Strictly impossible
under oracle correctness, suggesting KIT rounds arrival times down
via int16 deci-second quantization (the S2 effect) or operates on a
pre-truncated bundle.
\end{itemize}
The dominant bucket S5a (44.6\,pp) is not surgically attributable to
S2 alone without a KIT-internal patch; we therefore report it as
\emph{aggregate scan-side disagreement after subtracting the cleanly
attributable S1 cliff and S3 unreachability}. The full per-OD
classification CSV is at
\texttt{Experiments/results/cycle\_9\_paper/kit\_per\_suspect\_attribution/}\allowbreak\texttt{attribution\_3000.csv};
methodology and source in
\texttt{Experiments/results/cycle\_9\_paper/kit\_per\_suspect\_attribution/}\allowbreak\texttt{phase\_c\_attribute.py}.

\paragraph{S4 sanity check (multipass-blind cross-validation).}
Since both KIT and our oracle use single-pass CSA, the 185 multipass-blind
queries from Appendix~A should AGREE between KIT and oracle (both blind to
the same fixpoint relaxation). We observe: of those 185 q\_ids,
17 (9.2\%) AGREE, 4 NEAR, 137 DISAGREE, 23 ORACLE\_FINDS\_KIT\_BLIND,
4 DROPPED\_KIT --- the high DISAGREE share is dominated by S1/S3
suspects that are independent of the multipass blind spot, confirming
S4 does not provide a meaningful classification signal here.

\paragraph{Implication for the rest of the paper.} The benchmark
comparisons in main paper~\S6 use a \emph{within-pipeline} oracle that
is consistent across all evaluated methods (CSA, RAPTOR, ULTRA, TEEG,
MG, HL, CH). Hence the oracle-vs-KIT delta does not propagate to those
win-rate or speedup results; the KIT comparison in
Table~\ref{tab:kit-4way} is reported purely as an external
cross-implementation witness, not as ground truth for §6's tables.

For the static-latency comparison, the gap between KIT-CSA
(2.85\,ms/q) and our compiled CSA (12.74\,ms/q) reflects (a) C++ vs
Numba JIT and (b) KIT's SoA parallel-arrays layout vs our AoS
structured-numpy.

\section{Cycle 10 full matrix and post-update correctness oracle}
\label{app:cycle10}

\begin{table}[H]
\centering
\caption{Cycle 10: full 36-cell update timing matrix (3 cities $\times$
3 scenarios $\times$ 4 batch sizes). ``Update ms'' is wall-clock
median over the $N$ edits (single-threaded, warm cache, single run).
``Speedup'' denominator is the wall-clock time to rebuild the TEEG
data structure from scratch on the post-update timetable (London
571.3\,s, Berlin 138.5\,s, NYC 482.7\,s; see main paper Table~4).
This is the apples-to-apples cost a dynamic deployment would pay
without TEEG-INCR; it is NOT a comparison against a CSA query.
``Match'' $=$ (\# queries with post-update TEEG-INCR EA equal to
TEEG \emph{pre-update} baseline EA) / 3{,}000; this is a perturbation
measure. See Table~\ref{tab:rebuild-oracle} for the structural
correctness oracle against post-update CSA-rebuild.}
\label{tab:cycle10-full}
\scriptsize
\begin{tabular}{llrrrr}
\toprule
City & Op & $N$ & Update ms & Speedup & Match \\
\midrule
London & CANCEL & 1    & 12.5  & 45755$\times$ & 100.0\% \\
London & CANCEL & 10   & 14.8  & 38696$\times$ & 100.0\% \\
London & CANCEL & 100  & 22.8  & 25074$\times$ & 100.0\% \\
London & CANCEL & 1000 & 49.8  & 11460$\times$ & 96.0\% \\
London & DELAY  & 1    & 14.0  & 40925$\times$ & 100.0\% \\
London & DELAY  & 10   & 18.5  & 30862$\times$ & 99.8\% \\
London & DELAY  & 100  & 67.0  & 8520$\times$  & 99.8\% \\
London & DELAY  & 1000 & 389.9 & 1465$\times$  & 96.4\% \\
London & ADD    & 1    & 13.0  & 43952$\times$ & 100.0\% \\
London & ADD    & 10   & 16.9  & 33882$\times$ & 100.0\% \\
London & ADD    & 100  & 53.1  & 10757$\times$ & 100.0\% \\
London & ADD    & 1000 & 330.9 & 1726$\times$  & 99.4\% \\
Berlin & CANCEL & 1    & 5.0   & 27482$\times$ & 100.0\% \\
Berlin & CANCEL & 10   & 5.2   & 26469$\times$ & 100.0\% \\
Berlin & CANCEL & 100  & 8.2   & 16987$\times$ & 99.0\% \\
Berlin & CANCEL & 1000 & 17.7  & 7839$\times$  & 96.7\% \\
Berlin & DELAY  & 1    & 4.2   & 33149$\times$ & 100.0\% \\
Berlin & DELAY  & 10   & 6.0   & 22942$\times$ & 99.9\% \\
Berlin & DELAY  & 100  & 23.4  & 5920$\times$  & 99.7\% \\
Berlin & DELAY  & 1000 & 164.0 & 844$\times$   & 96.9\% \\
Berlin & ADD    & 1    & 4.9   & 28474$\times$ & 100.0\% \\
Berlin & ADD    & 10   & 6.1   & 22689$\times$ & 100.0\% \\
Berlin & ADD    & 100  & 17.2  & 8062$\times$  & 100.0\% \\
Berlin & ADD    & 1000 & 122.8 & 1128$\times$  & 99.9\% \\
NYC    & CANCEL & 1    & 7.7   & 62831$\times$ & 100.0\% \\
NYC    & CANCEL & 10   & 9.4   & 51498$\times$ & 100.0\% \\
NYC    & CANCEL & 100  & 19.9  & 24276$\times$ & 99.7\% \\
NYC    & CANCEL & 1000 & 46.8  & 10311$\times$ & 96.7\% \\
NYC    & DELAY  & 1    & 9.4   & 51610$\times$ & 100.0\% \\
NYC    & DELAY  & 10   & 15.3  & 31540$\times$ & 100.0\% \\
NYC    & DELAY  & 100  & 52.9  & 9119$\times$  & 99.8\% \\
NYC    & DELAY  & 1000 & 384.8 & 1255$\times$  & 95.9\% \\
NYC    & ADD    & 1    & 8.1   & 59444$\times$ & 100.0\% \\
NYC    & ADD    & 10   & 11.9  & 40477$\times$ & 100.0\% \\
NYC    & ADD    & 100  & 50.0  & 9651$\times$  & 100.0\% \\
NYC    & ADD    & 1000 & 376.8 & 1281$\times$  & 100.0\% \\
\bottomrule
\end{tabular}
\end{table}

\paragraph{Post-update CSA-rebuild oracle (3 cities, correctness validation).}

\begin{table}[H]
\centering
\caption{CSA-rebuild oracle on 3 cities $\times$ 12 cells $\times$ 300 OD
pairs = 10{,}800 evaluations. ``Match'' is the fraction of queries where
TEEG-INCR's post-update EA equals CSA-rebuild's. ``TB'' (teeg\_better)
counts cases where TEEG-INCR finds a strictly earlier EA (the
``17-residual'' phenomenon of Proposition~1). ``CB'' (csa\_better)
counts cases where CSA-rebuild finds a strictly earlier EA. Max
$|\Delta| \leq 15$\,min in every cell of every city.}
\label{tab:rebuild-oracle}
\scriptsize
\begin{tabular}{lr|rrrr|rrrr|rrrr}
\toprule
& & \multicolumn{4}{c|}{London} & \multicolumn{4}{c|}{Berlin} & \multicolumn{4}{c}{NYC} \\
Scen.\ & $N$ & M\% & TB & CB & $\Delta$ & M\% & TB & CB & $\Delta$ & M\% & TB & CB & $\Delta$ \\
\midrule
DELAY  & 1    & 99.7 & 1 & \textbf{0} & 12 & 99.7 & 1 & \textbf{0} & 2 & 99.3 & 2 & \textbf{0} & 9 \\
DELAY  & 10   & 99.7 & 1 & \textbf{0} & 12 & 99.7 & 1 & \textbf{0} & 2 & 99.3 & 2 & \textbf{0} & 9 \\
DELAY  & 100  & 99.3 & 1 & 1 & 12 & 99.7 & 1 & \textbf{0} & 2 & 99.3 & 2 & \textbf{0} & 9 \\
DELAY  & 1000 & 97.7 & 2 & 5 & 15 & 97.0 & 1 & 8 & 7 & 96.7 & 2 & 8 & 9 \\
CANCEL & 1    & 99.7 & 1 & \textbf{0} & 12 & 99.7 & 1 & \textbf{0} & 2 & 99.3 & 2 & \textbf{0} & 9 \\
CANCEL & 10   & 99.7 & 1 & \textbf{0} & 12 & 99.7 & 1 & \textbf{0} & 2 & 99.3 & 2 & \textbf{0} & 9 \\
CANCEL & 100  & 99.7 & 1 & \textbf{0} & 12 & 99.7 & 1 & \textbf{0} & 2 & 99.3 & 2 & \textbf{0} & 9 \\
CANCEL & 1000 & 99.3 & 1 & 1 & 12 & 99.3 & 2 & \textbf{0} & 6 & 99.7 & 1 & \textbf{0} & 3 \\
ADD    & 1    & 99.7 & 1 & \textbf{0} & 12 & 99.7 & 1 & \textbf{0} & 2 & 99.3 & 2 & \textbf{0} & 9 \\
ADD    & 10   & 99.7 & 1 & \textbf{0} & 12 & 99.7 & 1 & \textbf{0} & 2 & 99.3 & 2 & \textbf{0} & 9 \\
ADD    & 100  & 99.7 & 1 & \textbf{0} & 13 & 99.7 & 1 & \textbf{0} & 2 & 99.0 & 2 & 1 & 9 \\
ADD    & 1000 & 98.0 & 1 & 5 & 13 & 99.7 & 1 & \textbf{0} & 2 & 97.3 & 1 & 7 & 10 \\
\bottomrule
\end{tabular}
\end{table}

\paragraph{Cross-city aggregate.}
\begin{itemize}
\item \textbf{Total CB rate}: 36 / 10{,}800 = $0.33\%$
(London 0.33\%, Berlin 0.22\%, NYC 0.44\%); pattern consistent across cities.
\item \textbf{CB at $N \leq 10$ (realistic GTFS-RT cadence)}:
\textbf{0 / 5{,}400} across all 3 cities --- TEEG-INCR is \emph{exact}
in this regime.
\item \textbf{CB at $N\!=\!1000$ (stress test)}: 34 / 2700 = $1.26\%$
across all 3 cities (worst single cell: 8/300 = $2.67\%$ on
Berlin/NYC delay-$N\!=\!1000$).
\item \textbf{Max $|\Delta|$}: 15\,min on London, 10\,min on NYC,
7\,min on Berlin.
\end{itemize}
The pattern (TEEG-INCR exact for $N\!\leq\!10$, bounded
sub-millisecond-percent shortfall at $N\!=\!1000$) is robust across
all 3 cities; the shortfall is a deliberate property of the overlay's
single-link \textsc{board} construction (\S5.1) and not a function of
city topology. Per-query CSVs at
\texttt{cycle10\_rebuild\_oracle/\{,berlin\_,nyc\_\}oracle\_results.csv}.

\paragraph{Observations.}
For typical GTFS-Realtime cadences ($N \leq 10$ edits per query batch),
TEEG-INCR achieves \emph{exact} oracle agreement on $> 99.7\%$ of
queries on all three cities. At the stress-test $N\!=\!1000$, $\approx 1.26\%$
of queries show CSA-rebuild finding a strictly earlier EA
($\max|\Delta| \leq 15$\,min), traced to the single-link \textsc{board}
construction in the overlay: cascading cancellations create gaps in the
per-stop departure index that limit the overlay's transfer-path
coverage. This is a deliberate speed/completeness trade-off inherent
to $O(k)$ incremental patching; applications requiring exact
optimality under large batch updates can amortise a full rebuild
every $\sim\!1{,}000$ edits.

\paragraph{Per-cell classification of csa\_better cases.}
The 12 \emph{csa\_better} cases (out of 3{,}600 rebuild-oracle queries,
$0.33\%$) localise to the highest-stress cells with predominantly
$N \in \{100, 1000\}$ batch sizes:
\begin{itemize}
\item \texttt{delay\_N1000\_mag5}: 5 queries, $|\Delta|$ range
$\{1, 3, 4, 5, 10\}$ min
\item \texttt{add\_N1000\_mag5}: 5 queries, $|\Delta|$ range
$\{1, 2, 2, 9, 13\}$ min
\item \texttt{cancel\_N1000\_mag0}: 1 query, $|\Delta|=8$ min
\item \texttt{delay\_N100\_mag5}: 1 query, $|\Delta|=4$ min
\end{itemize}
No \emph{csa\_better} cases occur at $N \leq 10$ on any scenario,
confirming exactness in the typical GTFS-RT cadence. Per-query records
are in \texttt{cycle10\_rebuild\_oracle/oracle\_results.csv} (filter
\texttt{verdict == "csa\_better"}).

\section{NYC routes-array layout fix}
\label{app:nyc-fix}

The NYC GTFS feed was pre-processed with a different routes-array
column convention than London/Berlin
(\texttt{[rs\_lo,rs\_hi,st\_lo,st\_hi]} offset-pairs versus
\texttt{[n\_trips, n\_stops, rs\_lo, st\_lo]} count-+-offset). RAPTOR
and ULTRA read columns by index, so the wrong convention silently
returns junk indices into the stop-times CSR. The fix is a one-line
permutation:
$\text{n\_stops}=\text{rs\_hi}-\text{rs\_lo}$,
$\text{n\_trips}=(\text{st\_hi}-\text{st\_lo})/\text{n\_stops}$.
Code in \texttt{Cities/NYC/fix\_routes\_array\_layout.py}; validated
with $0/137{,}335$ unfound positions in
\texttt{build\_stop\_route\_position\_csr}. After fix: RAPTOR 8.12 ms
median (98.7\% coverage), ULTRA 8.24 ms median (98.7\% coverage); other
methods bit-identical on the 2{,}962 shared valid queries.

\section{MG-Dial-v2 $K\!=\!\infty$ fast path}
\label{app:mgdial-fix}

The original v2 query allocated 16 Dijkstra lanes regardless of the
query's transfer cap. For the unbounded-transfer case ($K\!=\!\infty$),
this wasted 16$\times$ memory and incurred lane-tracking overhead. The
fix adds a Numba kernel \texttt{\_dial\_single\_lane} that bypasses
lane logic and uses a single flat \texttt{int32[n\_nodes]} distance
array (mirroring v1), while preserving v2's tombstone-based dynamic
update primitives. After fix: v2 K=$\infty$ runs at $0.35\times$ v1
latency (\emph{faster} than v1 due to buffer reuse), 100\% correctness
on 20-query smoke test.

\section{Reproducibility manifest}
\label{app:repro}

The supplementary code release at
\url{https://anonymous.4open.science/r/teeg-incr-aaai2027-3CD9/}
(anonymous mirror; the underlying GitHub repository will be released
de-anonymised on acceptance) contains the artefacts needed to
regenerate every table and figure in the paper.

\paragraph{Data.}
GTFS feeds for the 2025-01-03 service window, per-city pre-processed
bundles \texttt{stop\_times\_df\_cyril.pkl}:
\begin{itemize}
\item London: source TfL GTFS bundle, snapshot 2025-01-02;
  \texttt{stop\_times\_df\_cyril.pkl} 1{,}279{,}201{,}489 bytes;
  SHA-256
  \texttt{cfd56dc60747ef27b59db1439c21b6c5\\bc9e736d4c8110e769870d97080f071e};
  raw \texttt{stop\_times} 2{,}946{,}672 rows
  $\to$ \texttt{\_collapse\_consecutive\_same\_stop} removes 565{,}143
  pairs $\to$ 2{,}381{,}529 rows $\to$ 07:00--20:00 window filter
  $\to$ 2{,}315{,}847 connections over 18{,}233 stops.
\item Berlin: source VBB GTFS bundle, snapshot 2025-01-02;
  \texttt{stop\_times\_df\_cyril.pkl} 229{,}847{,}646 bytes;
  SHA-256
  \texttt{9d8aba0fe820a374d55d365c9d95f925\\8a002852146832bc03c9decfbfc29200};
  10{,}816 stops, 781{,}877 connections.
\item NYC: source MTA GTFS bundle, snapshot 2025-01-02;
  \texttt{stop\_times\_df\_cyril.pkl} 421{,}789{,}907 bytes;
  SHA-256
  \texttt{a4b9989a9e504882f8c4269dafcb982f\\833f631f8497f25c16090e7bebb29ea0};
  14{,}908 stops, 1{,}394{,}454 connections.
\end{itemize}

\paragraph{Query sets.}
\texttt{queries\_3000\_\{london,berlin,nyc\}.csv}; sampler seed
\texttt{numpy.random.seed(20260101)} via
\texttt{Experiments/cycle\_runner/gen\_city\_queries.py}.

\paragraph{Scripts and the cells they produce.}
\begin{itemize}
\item Main Table~1 (per-city compiled bench):\\
  \texttt{Experiments/benchmark\_compiled\_only.py --queries Q --out O}\\
  produces \texttt{ms\_csa\_compiled}, \texttt{ms\_raptor\_compiled},
  \texttt{ms\_ultra\_compiled}, \texttt{ms\_teeg\_alt\_compiled},
  \texttt{ms\_ch\_compiled}, \texttt{ms\_hl\_compiled},
  \texttt{ms\_mg\_dial\_compiled} (the v1 columns for Table~1).\\
  \texttt{Experiments/benchmark\_compiled\_v2.py --queries Q --out O}\\
  produces \texttt{ms\_teeg\_alt\_v2}, \texttt{ms\_teeg\_alt\_incr},
  \texttt{ms\_mg\_dial\_v2} (the v2 columns for Table~1; median over
  rows 1..3000).\\
  CSV outputs:
  \texttt{p1a\_fullldn\_bench/bench\_3000\_fullldn\_compiled\_POSTFIX.csv},
  \texttt{p1a\_fullldn\_bench/bench\_3000\_fullldn\_v2.csv},
  \texttt{Cities/Berlin/bench\_3000\_berlin\_compiled.csv},
  \texttt{Cities/Berlin/bench\_3000\_berlin\_v2.csv},
  \texttt{Cities/NYC/bench\_3000\_nyc\_compiled\_FIXED.csv},
  \texttt{Cities/NYC/bench\_3000\_nyc\_v2.csv}. The Full London
  $K\!=\!\infty$ MG-Dial-v2 column (\texttt{ms\_mg\_dial\_v2} = 79.14
  median) is sourced from
  \texttt{mgdial\_v2\_rebench/bench\_fullldn\_postfix.csv}.
\item Cycle 10 dynamic-update matrix (Table~S2):\\
  \texttt{Experiments/cycle10\_run\_city.py --city CITY --queries Q}\\
  CSV: \texttt{Experiments/results/cycle\_10\_dynamic\_updates/RESULTS\_2026-06-09/\\cycle10\_results.csv}.
\item Cycle 10 rebuild-oracle correctness (Table~S3):\\
  \texttt{Experiments/results/cycle\_9\_paper/cycle10\_rebuild\_oracle/\\generate\_oracle.py}\\
  CSV: \texttt{cycle10\_rebuild\_oracle/oracle\_results.csv} (3{,}600 per-query rows).
\item KIT 4-way ablation (Table~S1):\\
  \texttt{Experiments/results/cycle\_9\_paper/kit\_4way\_ablation/\\scripts/\{build\_kit\_haversine\_walks, run\_cell3, compute\_agreement\}.py}.
\item Pareto figures (Figs 1--4):\\
  \texttt{Experiments/results/cycle\_9\_paper/paper\_figures/\\generate\_figures.py}.
\end{itemize}

\paragraph{Hardware.}
Bench machine: Intel Core i7-13700K (8P+8E cores, base 3.4 GHz),
64 GB DDR5-5600 RAM, Ubuntu 22.04 on WSL2, Python 3.12.4,
numpy 2.4.6, pandas 3.0.3, scipy 1.17.1, networkx 3.6.1, numba 0.65.1,
llvmlite 0.47.0. All bench runs single-thread with
\texttt{taskset -c 0} and CPU governor set to \texttt{performance}.

Cross-machine validation: Intel Xeon Platinum 8260L (8 vCPU,
2.30 GHz), 40 GB RAM, Samsung SSD 870 (369 MB/s), Ubuntu 22.04 LTS,
same Python+package versions; results within $\pm 18\%$ of dev box on
every cell of the lightweight \texttt{server\_validation\_package}
(KEY 1+2+3 of 4, with KEY 4 walkcap pending a re-run; full bundle
included in the supplementary release).

\paragraph{Repetition methodology.}
Main Table~1 values are medians over 3{,}000 queries from a single
bench run per (city, method) cell on machine A; $p_{95}$ and $p_{99}$
also recorded in the raw CSVs. Cycle 10 update timings are medians
over the $N$ edits in a single trace per cell.

\paragraph{3-rep Full London validation on machine B.}
To quantify run-to-run variance we re-ran the Full London compiled
bench (500 OD subset) three times on machine B (Intel Xeon Platinum
8260L, 8 vCPU, 40 GB RAM); per-method medians and across-rep spread
in Table~\ref{tab:multirep}. Per-method spread is $3\text{--}12\%$
across the 3 reps, well below the cross-method gaps that drive the
paper's conclusions.

\begin{table}[H]
\centering
\caption{Bench machine B (Xeon 8260L, 3 independent reps on Full
London 500-OD subset). ``Median'' is the median of 3 per-rep medians;
``Spread'' is $(\max - \min) / \text{median}$.}
\label{tab:multirep}
\small
\begin{tabular}{lrrrr}
\toprule
Method & Rep 0 & Rep 1 & Rep 2 & Spread \\
\midrule
HL              &   2.67 &   2.73 &   2.49 & 8.9\% \\
CSA             &   8.46 &   8.44 &   7.72 & 8.7\% \\
\textbf{TEEG-ALT-v2}    & \textbf{13.06} & \textbf{12.76} & \textbf{12.45} & \textbf{4.8\%} \\
TEEG-ALT-v1     &  13.44 &  13.14 &  12.38 & 8.0\% \\
TEEG-INCR       &  14.76 &  15.38 &  14.01 & 9.3\% \\
RAPTOR          &  16.75 &  16.81 &  14.91 & 11.4\% \\
ULTRA           &  17.32 &  17.32 &  16.01 & 7.6\% \\
CH              &  49.50 &  50.99 &  49.78 & 3.0\% \\
MG-Dij          & 190.76 & 198.37 & 175.26 & 12.1\% \\
MG-Dial         & 200.65 & 202.04 & 189.22 & 6.4\% \\
TEEG-B          & 280.66 & 269.70 & 251.34 & 10.9\% \\
\bottomrule
\end{tabular}
\end{table}

\paragraph{Cross-machine summary (post server-v2 validation 2026-06-13).}
\begin{itemize}
\item \emph{Berlin TEEG-ALT-v2 dominance over CSA}: machine
  A $2.34\times$, machine B (server pkg v2) $2.02\times$ ($1.88$ vs
  $3.79$ ms; 498 OD subset). Cross-machine dominance robust within $\pm 5\%$.
\item \emph{Cycle 10 Berlin oracle correctness}: machine B reproduces
  the dev-box correctness story --- 14/3600 CSA-better at N=1000
  (0.39\%), 36/3600 TEEG-better (1.00\%); 11 of 12 cells at
  $\geq 99\%$ match rate, single outlier cell add\_N1000 at 97.3\%
  match (server\_v2\_results/key7\_cycle10\_3city/berlin/summary.md).
\item \emph{17-residual multi-pass oracle}: machine B reproduces
  bit-for-bit --- multi-pass CSA matches TEEG-ALT on 16/17 residuals
  and strictly improves over single-pass on 17/17 (matches dev-box
  Appendix~A finding).
\item \emph{Full London TEEG-ALT-v2 vs CSA on machine B}:
  machine B with 3-rep bootstrap (498 OD subset) yields CSA $8.34$ ms
  CI95\% [$7.82, 8.69$] vs TEEG-ALT-v2 $12.69$ ms CI95\% [$11.57, 13.80$]
  with non-overlapping CIs (1.52$\times$ CSA-faster); machine A
  single-run reports the opposite ratio ($1.10\times$ TEEG-faster). We
  treat the Full London static-latency comparison as a within-hardware-class
  artefact and lead on Berlin/NYC for the dominance story.
\item \emph{Server pkg v2 partial-run note}: 3 of 7 KEYs completed
  successfully on the 40\,GB Xeon (KEY 3 cross-city, KEY 6 multi-pass,
  KEY 7 cycle 10 Berlin); 4 KEYs (KEY 1 Full LDN v2-vs-CSA, KEY 2
  Cycle 10 incr, KEY 4 walk-cap sweep, KEY 5 lazy-CSA) hit OOM
  (rc=$-9$) at the Full London 15-min walk-cap (12 GB TEEG +
  dual-residency overhead exceeded available memory). Server pkg v3
  (10-min walk-cap default + split-pass Cycle 10 + 6-point
  $\{3,5,7,10,12\,\text{min}\}$ walk-cap sweep) is in flight to
  recover those KEYs; see \texttt{Experiments/results/cycle\_9\_paper/}\allowbreak\texttt{server\_v2\_results/SUMMARY.md}.
\end{itemize}

\bibliographystyle{aaai2027}
\bibliography{references}

\end{document}

